\chapter*{A HYBRID KNOWLEDGE DISTILLATION AND PREFERENCE OPTIMIZATION FRAMEWORK FOR VIETNAMESE LEGAL REASONING IN SMALL LANGUAGE MODELS\\ABSTRACT}

Commercial Large Language Models are difficult to deploy in the legal domain due to data privacy, cost, and latency, while locally deployable Small Language Models lack multi-step legal reasoning. This thesis addresses this gap with a hybrid knowledge distillation and preference optimization framework that transfers reasoning ability from a large, reinforcement-learning-aligned teacher into lightweight Vietnamese legal students at the 0.6 and 1.7 billion parameter scales.

The framework comprises three phases organized into two tracks that share a supervised fine-tuned base. In the offline track, Offline Logit Distillation first trains the student on the teacher's soft target distributions under a supervised anchor and Quantized Low-Rank Adaptation; Diagnosis-Driven Preference Optimization then corrects residual errors on top of this checkpoint, using a Large Language Model as a Judge via Direct Preference Optimization. In parallel, the on-policy track applies On-Policy Knowledge Distillation directly to the same base, mitigating exposure bias by training on the student's own trajectories, guided by the teacher as an oracle.

Evaluated on three Vietnamese legal tasks, the 1.7B student after preference alignment achieves the strongest in-distribution reasoning performance on the NLI and syllogism tasks, quantized adaptation prevents catastrophic forgetting at the smallest scale, judge-based diagnosis yields cleaner preference data than rule-based extraction, and model scale proves decisive, as the advanced phases destabilize the 0.6B student. We further report a trade-off in which the supervised baseline generalizes better externally, demonstrating the feasibility of secure, compact legal reasoning assistants for on-premise deployment.
